% Generated draft: requires mathematical and source review.
\documentclass[11pt]{article}
\usepackage[margin=25mm]{geometry}
\usepackage{fontspec}
\setmainfont{Latin Modern Roman}
\usepackage{amsmath,amssymb,mathtools,unicode-math}
\setmathfont{Latin Modern Math}
\usepackage{enumitem,needspace,xurl,hyperref}
\setlength{\parindent}{0pt}
\setlength{\parskip}{5pt}
\setlength{\emergencystretch}{3em}
\newcommand{\N}{\mathbb{N}}
\newcommand{\Z}{\mathbb{Z}}
\newcommand{\Q}{\mathbb{Q}}
\newcommand{\R}{\mathbb{R}}
\newcommand{\C}{\mathbb{C}}
\newcommand{\F}{\mathbb{F}}
\newcommand{\A}{\mathbb{A}}
\newcommand{\PP}{\mathbb{P}}
\newcommand{\E}{\mathbb{E}}
\newcommand{\eps}{\varepsilon}
\newcommand{\dd}{\,\mathrm{d}}
\begin{document}
\section*{UnsolvedMath ID: 30000876}
\textbf{Source problem code:} catalogue entry
\subsection{Definitions and mathematical statement}
For a rational parametrization given by forms $f_0,\ldots,f_r$ in a parameter ring $B$, a moving relation of target degree $q$ is a polynomial homogeneous of degree $q$ in $X_0,\ldots,X_r$ that vanishes after substitution $X_i=f_i$. Its moving ideal can be modeled as the Rees-algebra kernel
\[
\ker\bigl(B[X_0,\ldots,X_r]\longrightarrow B[z],\quad X_i\longmapsto f_i z\bigr).
\]
Relations of target degree one are syzygies; higher target degrees give moving curves or surfaces in the intended parametrization setting.

Retain the source's three tasks:
\begin{enumerate}
\item Describe the ``remaining moving-curve ideal'' for rational parametrizations, including its proposed relation to Cox's conjecture and adjoint linear systems. The extract does not identify the already known part, the remaining generators, the version of Cox's conjecture or the adjoint system.
\item Describe the moving-surface ideal for tensor-product surfaces of bidegree $(2,1)$, including whatever generating or syzygy structure the intended family requires. The coordinate forms, map degree and base-point assumptions are not specified for this task.
\item For three bihomogeneous polynomials $f_1,f_2,f_3$ of bidegree $(2,1)$ with no common zero on $\PP^1\times\PP^1$, determine a free resolution of the ideal they generate.
\end{enumerate}
For the third task, make the standard parameter ring explicit:
\[
B=K[s_0,s_1,t_0,t_1],\qquad
\deg s_i=(1,0),\quad\deg t_j=(0,1),\qquad
I=(f_1,f_2,f_3)\subset B.
\]
A bigraded free resolution of $I$ begins with the presentation $B(-2,-1)^3\to I\to0$ and must specify the higher syzygy modules, shifts and maps; a minimal resolution is the usual refined target. The no-common-zero condition is on the projective product, not on affine four-space, and does not mean $I=B$. The question concerns the actual polynomial ideal, not merely its associated sheaf or a presumed complete-intersection quotient.

The ground field and geometric-zero convention, properness or birationality of parametrizations, and the exact moving-ideal and minimality conventions are not supplied. The three-generator ideal question is not silently identified with an entire surface-coordinate parametrization.
\subsection{Short English statement}
Describe the remaining moving-curve and bidegree-$(2,1)$ moving-surface ideals, and resolve the ideal of three base-point-free $(2,1)$ forms on $\PP^1\times\PP^1$. The connection to Cox's conjecture and adjoint systems is retained but not defined by the extract.
\subsection{Sources}
\begin{itemize}
\item Supplied report: \emph{Mini-Workshop: Surface Modeling and Syzygies} (2008), Oberwolfach Reports, DOI 10.4171/owr/2007/54, pp. 3181--3208; proposed year 2007. \url{https://doi.org/10.4171/owr/2007/54}.
\item Supplied reference labelled ``syzygies and the Rees algebra'': \url{https://www.sciencedirect.com/science/article/pii/S0022404907002897}.
\item Supplied reference labelled ``syzygy-based implicitization'': \url{https://arxiv.org/abs/1502.00890}.
\item Local UnsolvedMath registry, permanent ID 30000876, catalogue entry; created 25 August 2026, updated 16 September 2026. Its corrected\_verified note reports comparison with the official report. The supplied 21 August 2026 triage labels the combined record partially solved, reports important moving-ideal, syzygy and adjoint-system constructions, and does not claim a general resolution of all the remaining questions. This dated assessment is not independently updated here.
\end{itemize}
\end{document}
